% SOURCE_STATEMENT: "Let a Markov semigroup have a dilation to an automorphism or $E_0$-semigroup on an algebra $A$, and suppose $A$ is generated by the image of the original algebra and its time shifts. Must the dilation respect the GNS kernel, and must the induced $E_0$-semigroup on the represented subalgebra extend to all of $\\mathcal B^a(E)$? What weaker minimality hypotheses on $A$ suffice?"
% Preparation and final self-audit: CONTRIBUTING.md prompts C and D; source/audit artifacts retained in the local physics_statements_2026_10_11 batch.
% !TeX program = lualatex
\documentclass[11pt]{article}
\usepackage[margin=25mm]{geometry}
\usepackage{fontspec}
\setmainfont{Latin Modern Roman}
\usepackage{amsmath,amssymb,mathtools}
\usepackage{unicode-math}
\setmathfont{Latin Modern Math}
\usepackage{xurl,hyperref}
\hypersetup{colorlinks=true,urlcolor=blue}
\setcounter{secnumdepth}{0}
\setlength{\parindent}{0pt}
\setlength{\parskip}{5pt}
\setlength{\emergencystretch}{3em}
\newcommand{\sref}[2]{\hyperlink{source:#1:#2}{[S#2]}}
\newcommand{\cataloguescope}[1]{\paragraph{Recorded target.}#1}
\newcommand{\R}{\mathbb R}
\newcommand{\C}{\mathbb C}
\newcommand{\N}{\mathbb N}
\newcommand{\Z}{\mathbb Z}
\newcommand{\Q}{\mathbb Q}
\newcommand{\F}{\mathbb F}
\newcommand{\eps}{\varepsilon}
\begin{document}
% UnsolvedMath ID 30001184; source catalogue code OWR-3392-012
\setcounter{section}{30001183}
\section{Extending Dilated E0-Semigroups Beyond GNS Representations}\label{30001184}
{\small UnsolvedMath ID 30001184 · Mathematical Physics}
\subsection{Definitions and mathematical statement}

\paragraph{Shared notation.}$\mathbb N=\{1,2,\ldots\}$; the standard number fields have their usual meanings. All additional conventions are stated locally.


\paragraph{Definitions and conventions.} A Markov semigroup $T_t$ is a continuous semigroup of unital completely positive maps on the specified operator algebra $B$. A dilation consists of an embedding $i:B\to A$, a conditional expectation or compression $p:A\to B$, and endomorphisms $\vartheta_t$ with $T_t=p\vartheta_t i$. The GNS construction for $p$ gives a Hilbert $B$-module $E$ and a represented algebra in $\mathcal B^a(E)$. An automorphic dilation has invertible $*$-automorphisms; an inner one is implemented by unitaries. Respecting the GNS kernel means that the kernel of this representation is carried into itself, so the dilation descends to the represented algebra. Extending to $\mathcal B^a(E)$ means finding unital endomorphisms of the full adjointable-operator algebra agreeing with that descended semigroup.

\paragraph{Source scope and qualifications.} Both kernel invariance and the full-algebra extension are needed. The phrase not too big is conjectural; generation by i(B) and its time shifts is the report's proposed example of minimality, not a proved sufficient condition.

\paragraph{Catalogue mathematical target.}
\paragraph{Catalogue wording retained for traceability.}
Let a Markov semigroup have a dilation to an automorphism or $E_0$-semigroup on an algebra $A$, and suppose $A$ is generated by the image of the original algebra and its time shifts. Must the dilation respect the GNS kernel, and must the induced $E_0$-semigroup on the represented subalgebra extend to all of $\mathcal B^a(E)$? What weaker minimality hypotheses on $A$ suffice?

\paragraph{Saved literature context.} The supplied catalogue's assessment dated 2026-08-21 records: The report states this as a conjecture about extending the induced E0-semigroup and explicitly asks for a proof or counterexample; no general resolution was located. The assessment is retained as dated provenance; it is not a new comprehensive literature or open-status certification.
\subsection{Short English statement}
Find minimality assumptions ensuring that a general Markov dilation descends through its GNS representation and extends to all adjointable module operators.
\subsection{Sources}
\begin{itemize}
\item[\hypertarget{source:30001184:1}{S1}] ULAM.AI and the UnsolvedMath Contributors, UnsolvedMath, record 30001184 (OWR-3392-012). Catalogue curation under CC BY 4.0; definitions and English exposition added locally.
\url{https://huggingface.co/datasets/ulamai/UnsolvedMath}.
\item[\hypertarget{source:30001184:2}{S2}] Mini-Workshop: Product Systems and Independence in Quantum Dynamics (2009), Oberwolfach Reports, DOI 10.4171/owr/2009/09, pp. 493--548; official report source and mathematical context.
\url{https://doi.org/10.4171/owr/2009/09}.
\item[\hypertarget{source:30001184:3}{S3}] Official EMS Press PDF of the same Oberwolfach report; PDF pages 50--51, printed pages 542--543.
\url{https://ems.press/content/serial-article-files/46209}.
\end{itemize}
\label{end:30001184}
\end{document}
